\subsection{Photon-Photon Interactions}
As mentioned above, photons do not interact with each other in first order. Their interaction comes about through virtual fermions in higher-order loop diagrams within the framework of quantum electrodynamics (QED). In QED, virtual particle-antiparticle pairs such as electron-positron pairs can be created from the vacuum due to excitations of the fermion fields by the photon. These virtual loops allow for a non-trivial interaction between two real photons, effectively enabling a photon–photon scattering process. This interaction arises through a second-order Feynman diagram where one photon emits a virtual electron-positron pair, which is then reabsorbed by the second photon.

%!------------------Add illustration of photon-photon interaction through loop-----------------!

\subsection{Photon-Hadron Interactions}
In addition, theory suggests real (or quasi-real) photons can also interact with hadrons in two different ways: Through Vector Meson Dominance (VMD) fluctuations or through QCD fluctuations to $q\bar{q}$ pairs. To discuss such descriptions, Schuler and Sjöstrand \cite{schuler_towards_1993} describe the wave function of the photon as a superposition:

\begin{equation}\label{eq:photon_wavefunction}
	\ket{\gamma} = \sqrt{Z_3}\ket{\gamma_{B}} + \sum\limits_{V=\rho^0,\omega,\phi}\frac{e}{f_{V}}\ket{V} + \frac{e}{f_{q\bar{q}}}\ket{q\bar{q}} \,,
\end{equation}

where $Z_3$ is $1-\Sigma_{V=\rho^0,\omega,\phi}\left(\frac{e}{f_{V}}\right)^2 - \left(\frac{e}{f_{q\bar{q}}}\right)^2$, $\ket{\gamma_B}$ is the ``bare'' photon, or the photon component as point particle we associate with EM theory and QED, and therefore cannot interact at all with hadrons. $\ket{V}$ is the vector meson component and $\ket{q\bar{q}}$ is the hadron component. The summation on $\rho^0,\omega,\phi$ is due to them being the only hadronic constituents of the photon under simple VMD. VMD description however, breaks down in small $Q^2$ and small $x$, at which the final expression of the photon $\ket{q\bar{q}}$ takes effect and become the contributor to photon-parton interactions. In the proposal, we shall refer to photon interactions in which the latter is in effect as ``resolved'' photon interaction, to denote the dissociation into $q\bar{q}$ pair which will then interact with the target.


\begin{figure}[h!] 
	\centering % Centers the image and caption
	\includegraphics[width=0.8\textwidth]{\main/illustrations/schuler_photn_proton_interactions.png}
	\caption{Contributions to hard $\gamma$p interactions: (a) direct, (b) VMD, and (c) anomalous. Only the basic graphs are illustrated; additional partonic activity is allowed in either case. The presence of spectator jets has been indicated by dashed lines, while full lines show partons that (may) give rise to high-$p_T$ jets \cite{schuler_towards_1993}.}
	\label{fig:schuler_photon_proton_diagrams}
\end{figure}


\subsection{Photon-Pomeron Interactions}

As mentioned in \ref{sec:pomeron}, the pomeron photon interactions are a form of effective theory to describe collisions in which the photon interacts with a charge, color neutral complex QCD object during diffraction events. An example for such a process can be seen in Fig. \ref{fig:gammaP_interaction}. Here, as described above, the photon contribution is through a resolved state by dissociation into a $q\bar{q}$ pair, which then interacts with the pomeron by means of gluon exchange. It is important to notice that through this exchange there will be additional particle production from the pomeron, denoted in X.

\begin{figure}[h!] 
	\centering % Centers the image and caption
	\includegraphics[width=0.8\textwidth]{\main/illustrations/gammaP_interaction.png}
	\caption{Typical leading-order Feynman graphs for diffractive di-jet photoproduction in UPCs of hadrons A and B. Graphs a and b correspond to the direct and resolved photon contributions,respectively \cite{guzey_diffractive_2016}.}
	\label{fig:gammaP_interaction}
\end{figure}

 

There are two distinct types of photon-pomeron interaction, which are often called coherent and incoherent interactions and stem from the pomeron itself. In coherent interaction, the pomeron mediates an exchange from the entire nucleus, while in incoherent interaction the pomeron is associated with a single nucleon. Incoherent interaction then leaves the probed nucleus either in an excited state or causing it to break up. This is in contrast to coherent interaction in which the nucleus will most likely remain intact.  

